\section*{Chapter \ref{ch:m1:matching-algorithms-and-implied-spreads} --- Matching Algorithms and Implied Spreads}

\begin{solution}{exo:m1:matching-algorithms-and-implied-spreads:1}
A 10, B 200, C 20 of its 40, D nothing.
\end{solution}

\begin{solution}{exo:m1:matching-algorithms-and-implied-spreads:2}
Shares $60 \times (10, 200, 40, 250)/500 = 1.2$, 24, 4.8, 30. Floors 1, 24, 4,
30: 59 lots. The one lot left over goes to the oldest order, A. Fills: A 2, B
24, C 4, D 30.
\end{solution}

\begin{solution}{exo:m1:matching-algorithms-and-implied-spreads:3}
Bid $10\,050 - 10\,023 = 27$ for $\min(30, 10) = 10$. Offer $10\,052 - 10\,020 =
32$ for $\min(25, 40) = 25$.
\end{solution}

\begin{solution}{exo:m1:matching-algorithms-and-implied-spreads:4}
Reducing the size keeps the place. Increasing it, changing the price, and
cancelling and re-entering all lose it (the last by construction). To add
size without losing priority, never amend upwards: leave the old order where
it is and enter a \emph{second} order for the additional quantity, which
joins the back of the queue. The order manager must therefore handle several
live orders per price level and cancel the youngest first when reducing.
\end{solution}

\begin{solution}{exo:m1:matching-algorithms-and-implied-spreads:5}
$8 \times 1\,000/92 = 86.96$, so 87; $\lfloor 100 \times 87/1\,087\rfloor = 8$,
$\lfloor 100 \times 86/1\,086\rfloor = 7$, $\lfloor 1\,000 \times 87/1\,087\rfloor
= 80$. If everyone multiplies by ten, shares are unchanged: each receives
what it received before, the displayed depth is ten times larger, and each is
exposed to ten times more on a sweep. Displayed depth then measures the
participants' tolerance for being over-filled, not their wish to trade, which
is why depth in pro-rata products evaporates in fast markets.
\end{solution}

\begin{solution}{exo:m1:matching-algorithms-and-implied-spreads:6}
One order of 15 lots was cancelled (a trade would have produced a trade
message). If it was ahead of yours, 15 fewer lots are ahead; if behind,
nothing changed. Market by price does not say which. A common estimate
assumes cancellations are spread in proportion to the quantity ahead and
behind, or, more conservatively, that they all come from behind. Market by
order gives the cancelled order's identifier and hence its place: the
position is known exactly.
\end{solution}

\begin{solution}{exo:m1:matching-algorithms-and-implied-spreads:7}
Front of a FIFO queue: 9.56 of the 10 lots. Pro rata: 2.34 wherever the order
stands. Ratio 4.1. Pro rata is better for an order with 170 lots or more
ahead of it in a 500-lot level.
\end{solution}

\begin{solution}{exo:m1:matching-algorithms-and-implied-spreads:8}
The exact share of a 100-lot aggressor is $100 \times 5/2\,005 = 0.249$ lot: it
rounds down to zero, and with a two-lot minimum nothing smaller than 41
lots shown ($\lfloor 100 \times 41/2\,041\rfloor = 2$) is allocated anything
pro rata. Leftovers go by time to older orders. The trader collects nothing
from ordinary flow and is filled only when an order of 800 lots or more
sweeps the level: a portfolio of exclusively bad fills. In a pro-rata product
there is a minimum viable size, and it is large.
\end{solution}

\begin{solution}{pb:m1:matching-algorithms-and-implied-spreads:1}
\textbf{1.} $10\,050 - 10\,023 = 27$ for 10.
\textbf{2.} $10\,052 - 10\,020 = 32$ for 25.
\textbf{3.} Bid 28 for 5 (direct); ask 32 for 32 (7 direct and 25 implied).
\textbf{4.} It would mean the front bid exceeds the back ask plus a price
at which someone is willing to sell the spread: the engine would match the
three orders at once. Crossed implieds against direct orders trade
immediately; they are never displayed.
\textbf{5.} An implied bid and an implied ask that crossed would consist
entirely of outright orders that do not cross in their own books; matching
them would require chains of legs in further instruments, without limit.
Exchanges cut the recursion: every match contains at least one direct
order.
\textbf{6.} The direct bid, 5 at 28; then the implied bid, 7 at 27.
\textbf{7.} $(5 \times 28 + 7 \times 27)/12 = 27.42$.
\textbf{8.} For the implied part: 7 Z bought by the front-month bidder at
10\,050 and 7 H sold by the back-month seller at 10\,023. The direct part
prints in the spread only (the legs are assigned prices by an exchange
convention and do not touch the outright books).
\textbf{9.} Front: bid 10\,050 $\times$ 23, ask unchanged. Back: bid
unchanged, ask 10\,023 $\times$ 3. Spread: bid 27 $\times$ 3 (implied), ask 32
$\times$ 32.
\textbf{10.} Sell 12 Z at 10\,050; buy 10 H at 10\,023 and 2 at the next
offer, 10\,024 at best: a spread of at most 26.83, against 27.42, with the
risk of the market moving between the two legs.
\textbf{11.} Bid $28 + 10\,020 = 10\,048$ for 5; ask $32 + 10\,023 = 10\,055$ for
7.
\textbf{12.} No: the front month shows 10\,050 at 10\,052, better on both
sides.
\textbf{13.} $10\,050 - 10\,021 = 29$ for $\min(30, 20) = 20$.
\textbf{14.} No. In its own book the best bid is 10\,020, and the \omterm{def:m1:matching-algorithms-and-implied-spreads:implied}{implied-out}
back bid, $b_F - a_S = 10\,050 - 32 = 10\,018$, is lower still. It would trade
if someone sold the spread at 29 or better, or bid 10\,021 in the back month.
\textbf{15.} One tick behind: the displayed best bid is now 29 for 20,
implied. To be at the front again the \omterm{def:m1:what-a-trading-firm-does:market-maker}{market maker} must bid 29, and where
that order ranks against the implied 20 lots depends on the engine's rule.
An order in the back month has just repriced the spread without any spread
order being entered: quotes in one book must be driven by all three.
\textbf{16.} It misses fills that come from spread sellers through
\omterm{def:m1:matching-algorithms-and-implied-spreads:implied}{implied-out} orders, and it misjudges queue position: \omterm{def:m1:matching-algorithms-and-implied-spreads:implied}{implied-out} quantity
can stand ahead of, or behind, the strategy's order at the same price. It
also misses that visible front-month depth includes quantity which vanishes
when the \emph{back} month moves.
\textbf{17.} If implieds rank behind all direct orders, a direct order's
queue position is safe; if they rank by the time they were created, implied
quantity can step ahead whenever another book moves. The expected fill, and
its \omterm{def:m1:what-a-trading-firm-does:adverse-selection}{adverse selection}, differ between the two.
\textbf{18.} Computing and disseminating implied depth at every level of
every combination is expensive and changes with every update in any leg;
the best level is what matters for matching.
\textbf{19.} \textbf{Implied best bid 27 for 10; displayed best bid 28 for 5.}
\textbf{20.} It lets a roll entered as one order draw on the liquidity of
both expiries at once, without the risk of being filled on one leg only.
\end{solution}

\begin{solution}{iq:m1:matching-algorithms-and-implied-spreads:1}
FIFO fills resting orders at the best price oldest first; pro rata shares
the incoming order in proportion to size. Actively moving products with
small ticks relative to volatility (equity-index, crude, note futures) are
FIFO: priority is earned by improving the price or by arriving first.
Products that barely move relative to their tick, short-term interest-rate
futures above all, use pro rata or hybrids with top-order and
market-maker allocations, so that many firms keep quoting at a price that
stays put for hours.
\iqlookfor{the link between tick size relative to volatility and the rule.}
\end{solution}

\begin{solution}{iq:m1:matching-algorithms-and-implied-spreads:2}
An order the matching engine derives from orders in related books: two
outright orders in different expiries imply a calendar-spread order
(\omterm{def:m1:matching-algorithms-and-implied-spreads:implied}{implied-in}); a spread order and an outright imply an order in the other leg
(\omterm{def:m1:matching-algorithms-and-implied-spreads:implied}{implied-out}). It is tradable, all legs execute atomically, and implied
never matches implied.
\iqlookfor{both directions, atomic execution, no implied against implied.}
\end{solution}

\begin{solution}{iq:m1:matching-algorithms-and-implied-spreads:3}
FIFO: invest in latency to reach new price levels first, keep orders alive to
preserve queue position, value each position in the queue, show true size.
Pro rata: speed matters less; show more than intended to earn a share,
manage the risk of being over-filled on sweeps, cancel quickly when large
orders appear, and care about minimum allocation thresholds and any
top-order privilege, which brings the speed race back for the first order at
a new price.
\iqlookfor{oversizing and its risk; the top-order step reintroducing speed.}
\end{solution}

\begin{solution}{iq:m1:matching-algorithms-and-implied-spreads:4}
Compute $\lfloor a \cdot s_i / S\rfloor$ in integers for each order, on open
quantities; drop allocations below the minimum; distribute the remainder by
the secondary rule. Edge cases: aggressor larger than the level; total of
floors less than the aggressor; all allocations below the minimum (then
everything goes by the secondary rule); an order already partly filled by an
earlier step; overflow of $a \cdot s_i$; never exceeding an order's size;
determinism of the tie-break. Test conservation on random books.
\iqlookfor{integer arithmetic, conservation, and the all-below-minimum case.}
\end{solution}

\begin{solution}{iq:m1:matching-algorithms-and-implied-spreads:5}
Record the level's quantity when the order is acknowledged: that is the
quantity ahead. Trades at the level reduce it one for one. Decreases without
trades are cancellations: allocate them ahead and behind by an assumption
(proportional, or all behind for a conservative estimate), bounded so that
the quantity ahead never exceeds the level's quantity minus one's own order.
Calibrate the assumption against one's own fills: the realised fill times
reveal the true position after the fact.
\iqlookfor{the bound from the level's total, and calibration on own fills.}
\end{solution}

\begin{solution}{iq:m1:matching-algorithms-and-implied-spreads:6}
A single engine over all legs, since an update in any book can create or
remove tradable quantity in the others; atomic multi-leg execution with leg
prices assigned by the exchange's convention; the rule ranking implied
against direct orders at equal prices; no implied-against-implied; different
ticks in spreads and outrights; the allocation algorithm of each book; and
market data that reports implied quantity the way the real feed does,
otherwise strategies trained on the simulator will see a book that does not
exist.
\iqlookfor{cross-book event propagation and the ranking rule.}
\end{solution}
